% Generated by scripts/assemble.py; corrections belong in source/corrections.json.
\examhead{122}{山东大学 2025 年数学分析试卷}{117}
% p118-b01; source PDF page 118
\par\Needspace{5\baselineskip}
\blockanchor{p118-b01}
\majorquestion{一}计算题. 每题 15 分, 共 45 分.\par
% p118-b02; source PDF page 118
\par\Needspace{6\baselineskip}
\blockanchor{p118-b02}
\question{1}设 $g(x)$ 在 $[0,1]$ 上二阶连续可微, 满足\par
\nopagebreak[4]
% p118-b03; source PDF page 118
\blockanchor{p118-b03}
\displayformula{g (0) = g ^ {\prime} (0) = 0, g ^ {\prime \prime} (0) = 2 a.}
% p118-b04; source PDF page 118
\blockanchor{p118-b04}
设 $f(x) = \frac{g(x)}{x} (x\neq 0),f(0) = 0,$ 求 $f^{\prime}(0)$\par
% p118-b05; source PDF page 118
\blockanchor{p118-b05}
\question{2}求积分 $\int_0^1\frac{\ln(1 + x)}{1 + x^2}\mathrm{d}x.$\par
% p118-b06; source PDF page 118
\blockanchor{p118-b06}
\question{3}若 $F(t) = \iiint_{x^2 + y^2 + z^2 \leq t^2} f(x^2 + y^2 + z^2) \, \mathrm{d}x \, \mathrm{d}y \, \mathrm{d}z, t > 0$. 其中 $f$ 是 $\mathbb{R}$ 上的可微函数.\par
% p118-b07; source PDF page 118
\blockanchor{p118-b07}
(i) 计算  $F'(t)$ .\par
% p118-b08; source PDF page 118
\blockanchor{p118-b08}
(ii) 若 $f'(0) = 0$, 计算 $\lim_{t \to 0^+} \frac{F(t)}{t^5}$.\par
% p118-b09; source PDF page 118
\par\Needspace{5\baselineskip}
\blockanchor{p118-b09}
\majorquestion{二}证明题. 每题15分, 共45分.\par
% p118-b10; source PDF page 118
\blockanchor{p118-b10}
\question{1}设 $f(x)$ 在 (0,1) 上可导, $f'(x)$ 在 (0,1) 上有界, 证明: $f(x)$ 在 (0,1) 上有界且一致连续.\par
% p118-b11; source PDF page 118
\blockanchor{p118-b11}
\question{2}设 $f(x, y) = \varphi(|xy|)$，且在 $u = 0$ 的附近满足 $|\varphi(u)| \leq u^2$. 证明：$f(x, y)$ 在 (0,0) 处可微.\par
% p118-b12; source PDF page 118
\par\Needspace{6\baselineskip}
\blockanchor{p118-b12}
\question{3}设非负连续函数 $f(x)$ 在 $[0, +\infty)$ 上单调递减，证明：\par
\nopagebreak[4]
% p118-b13; source PDF page 118
\blockanchor{p118-b13}
\displayformula{F (x) = \int_ {0} ^ {x} (x - 2 t) f (t) \mathrm{d} t}
% p118-b14; source PDF page 118
\blockanchor{p118-b14}
是 $[0, +\infty)$ 上的凸函数.\par
% p118-b15; source PDF page 118
\par\Needspace{5\baselineskip}
\blockanchor{p118-b15}
\majorquestion{三}综合题. 每题 20 分, 共 60 分.\par
% p118-b16; source PDF page 118
\blockanchor{p118-b16}
\question{1}设 $f(x) = \frac{1}{1 - x^2}$.\par
% p118-b17; source PDF page 118
\blockanchor{p118-b17}
(1) 求极限 $\lim_{n\to \infty}\left[f\left(1 - \frac{1}{n + 1}\right) - f\left(1 - \frac{1}{n}\right)\right].$\par
% p118-b18; source PDF page 118
\blockanchor{p118-b18}
(2) 证明: $f(x)$ 在 $[0,1)$ 上不一致连续.\par
% p118-b19; source PDF page 118
\par\Needspace{6\baselineskip}
\blockanchor{p118-b19}
\question{2}设\par
\nopagebreak[4]
% p118-b20; source PDF page 118
\blockanchor{p118-b20}
\displayformula{u _ {n} (x) = x ^ {n} \ln x,   x \in (0, 1 ].}
% p118-b21; source PDF page 118
\blockanchor{p118-b21}
讨论 $\sum_{n=1}^{\infty} u_n(x)$ 在 $(0,1]$ 上的收敛性与一致收敛性, 并计算 $\int_0^1 \sum_{n=1}^{\infty} u_n(x) \, \mathrm{d}x$.\par
% p118-b22; source PDF page 118
\blockanchor{p118-b22}
\question{3}设 $u(x, y, z)$ 是 $\mathbb{R}^3$ 上的连续函数, 而且在点 $(x_0, y_0, z_0)$ 的某个邻域上有连续的二阶偏导数. 记 $\Sigma$ 是以点 $(x_0, y_0, z_0)$ 为球心, $R$ 为半径的球面, 令 $T(R) = \frac{1}{4\pi R^2} \iint_{\Sigma} u(x, y, z) \mathrm{d}S$, $R > 0$. 证明:\par
% p118-b23; source PDF page 118
\blockanchor{p118-b23}
(1) $\lim_{R\to 0^{+}}T(R) = u(x_0,y_0,z_0)$.\par
% p118-b24; source PDF page 118
\blockanchor{p118-b24}
(2) $\lim_{R\to 0^{+}}\frac{T(R) - u(x_0,y_0,z_0)}{R^2} = \frac{1}{6}\Delta u(x_0,y_0,z_0)$, 其中 $\Delta u = \frac{\partial^2u}{\partial x^2} +\frac{\partial^2u}{\partial y^2} +\frac{\partial^2u}{\partial z^2}.$\par
